\subsection{Structural Summary}

The preceding sections separate two ideas that are often confused:

\begin{itemize}
    \item Python names are dynamically bound.
    \item Python objects are strongly typed at runtime.
\end{itemize}

This means that Python is flexible at the level of assignment, but strict at the level of operation. A name may point to different objects at different moments, but each object still brings its own type rules with it.

\begin{verbatim}
value = 42
print(f"value={value!r:<10} type={type(value)}")

value = "42"
print(f"value={value!r:<10} type={type(value)}")

value = None
print(f"value={value!r:<10} type={type(value)}")
\end{verbatim}

Possible output:

\begin{verbatim}
value=42         type=<class 'int'>
value='42'       type=<class 'str'>
value=None       type=<class 'NoneType'>
\end{verbatim}

The name \texttt{value} moved. The objects did not lose their type identity.

\subsubsection{Categorizing Python as a Dynamically Bound, Strongly Verified Runtime Typing Environment}

Python is dynamically typed because names do not have fixed declared types.

There is no declaration such as:

\begin{verbatim}
int value;
\end{verbatim}

Instead, assignment binds a name to an object:

\begin{verbatim}
value = 42
\end{verbatim}

Later, the same name can be rebound:

\begin{verbatim}
value = "No soup for you."
\end{verbatim}

A compact trace:

\begin{verbatim}
value = 42

print("first binding")
print(f"value: {value!r}")
print(f"type:  {type(value)}")
print(f"id:    {id(value)}")

value = "No soup for you."

print()
print("second binding")
print(f"value: {value!r}")
print(f"type:  {type(value)}")
print(f"id:    {id(value)}")
\end{verbatim}

Possible output:

\begin{verbatim}
first binding
value: 42
type:  <class 'int'>
id:    140734432946520

second binding
value: 'No soup for you.'
type:  <class 'str'>
id:    2209266113904
\end{verbatim}

This is dynamic binding.

However, Python is also strongly checked at runtime. The interpreter does not ignore object types when operations execute.

\begin{verbatim}
left = "3"
right = 4

try:
    result = left + right
    print(result)
except TypeError as err:
    print("operation failed")
    print(f"message: {err}")
    print(f"type(err): {type(err)}")
\end{verbatim}

Possible output:

\begin{verbatim}
operation failed
message: can only concatenate str (not "int") to str
type(err): <class 'TypeError'>
\end{verbatim}

The failure is not caused by the name \texttt{left}. It is caused by the object reached through \texttt{left}. That object is a string, and string addition does not accept an integer as the right operand.

The correct operation must be stated explicitly:

\begin{verbatim}
text_number = "3"
real_number = 4

numeric_result = int(text_number) + real_number
text_result = text_number + str(real_number)

print(f"numeric_result: {numeric_result!r}, type={type(numeric_result)}")
print(f"text_result:    {text_result!r}, type={type(text_result)}")
\end{verbatim}

Possible output:

\begin{verbatim}
numeric_result: 7, type=<class 'int'>
text_result:    '34', type=<class 'str'>
\end{verbatim}

So the proper classification is not simply ``dynamic''.

A more precise description is:

\begin{quote}
Python uses dynamically bound names and strongly verified runtime operations.
\end{quote}

That means:

\begin{itemize}
    \item names can be rebound freely;
    \item objects carry runtime types;
    \item operators consult object behavior;
    \item unsupported combinations fail instead of being silently converted;
    \item explicit conversion is the programmer's responsibility.
\end{itemize}

\subsubsection{The Practical Rule: Names Are Flexible, Objects Keep Their Runtime Type}

The practical rule for reading Python code is:

\begin{quote}
A name may change what it references, but the object it currently references keeps its type and behavior.
\end{quote}

Example:

\begin{verbatim}
thing = 42

print("integer phase")
print(f"thing: {thing!r}")
print(f"type:  {type(thing)}")
print(f"bit_length available: {'bit_length' in dir(thing)}")
print(f"upper available:      {'upper' in dir(thing)}")

thing = "D'oh!"

print()
print("string phase")
print(f"thing: {thing!r}")
print(f"type:  {type(thing)}")
print(f"bit_length available: {'bit_length' in dir(thing)}")
print(f"upper available:      {'upper' in dir(thing)}")
\end{verbatim}

Possible output:

\begin{verbatim}
integer phase
thing: 42
type:  <class 'int'>
bit_length available: True
upper available:      False

string phase
thing: "D'oh!"
type:  <class 'str'>
bit_length available: False
upper available:      True
\end{verbatim}

The name \texttt{thing} did not carry the methods. The current object did.

The same rule explains operation behavior:

\begin{verbatim}
thing = 42
print(thing + 8)

thing = "D'oh!"
print(thing + "!")
\end{verbatim}

Possible output:

\begin{verbatim}
50
D'oh!
\end{verbatim}

Both uses of \texttt{+} work, but they do not mean exactly the same operation:

\begin{itemize}
    \item for integers, \texttt{+} performs numeric addition;
    \item for strings, \texttt{+} performs text concatenation.
\end{itemize}

The object type decides which behavior is available.

If the operation is not defined, Python rejects it:

\begin{verbatim}
thing = "42"

try:
    print(thing + 8)
except TypeError as err:
    print("operation failed")
    print(f"message: {err}")
\end{verbatim}

Possible output:

\begin{verbatim}
operation failed
message: can only concatenate str (not "int") to str
\end{verbatim}

A final compact script:

\begin{verbatim}
values = [42, 3.14, 2 + 5j, True, None, "No soup for you."]

for value in values:
    print(f"{repr(value):<20} type={type(value)}")
\end{verbatim}

Possible output:

\begin{verbatim}
42                   type=<class 'int'>
3.14                 type=<class 'float'>
(2+5j)               type=<class 'complex'>
True                 type=<class 'bool'>
None                 type=<class 'NoneType'>
'No soup for you.'   type=<class 'str'>
\end{verbatim}

Every value is an object with a runtime type. A name can point to any of them, but it does not erase their differences.

The structural summary for this whole section is:

\begin{verbatim}
assignment:
    flexible name-to-object binding

reassignment:
    same name, possibly different object type

type():
    exact current runtime type

isinstance():
    category or family test

operation:
    checked against the current objects involved

conversion:
    explicit when crossing type boundaries
\end{verbatim}

The final rule is:

\begin{quote}
Python is permissive about rebinding names, but strict about what objects can do together.
\end{quote}